\documentclass{article}
\usepackage[UTF8]{ctex}
\usepackage{amsmath,amssymb,amsfonts}
\usepackage{enumitem}
\usepackage{geometry}
\geometry{a4paper, margin=1in}
% === 错题本排版优化（不参与题目标识） ===
\setlength{\parindent}{0pt}
\setlength{\parskip}{0.35em plus 0.10em minus 0.08em}
\setlength{\emergencystretch}{3em}
\setlist{itemsep=0.25em, topsep=0.40em, parsep=0pt}
\allowdisplaybreaks[2]
% === 错题本排版优化结束 ===
\begin{document}
\section*{北京工业大学2025年数学分析真题}
\section*{9．北京工业大学2025年数学分析试卷}
\begin{itemize}

\subsection{题目1}
  \item[1．] （15 分）设 \(p>1\) ，证明：数列 \(\left\{1+\frac{1}{2^{p}}+\cdots+\frac{1}{n^{p}}\right\}\) 收敛，且其极限 \(a\) 满足
\[
1+\frac{1}{2^{p-1}(p-1)} \leq a \leq \frac{p}{p-1}
\]

\subsection{题目2}
  \item[2．] （15 分）证明：方程组 \(\frac{5}{x-1}+\frac{7}{x-2}+\frac{16}{x-3}=0\) 在 \((1,2)\) 与 \((2,3)\) 内各有一实根．
\end{itemize}

\subsection{题目3}
3．（15 分）设函数 \(f(x)\) 在 \([0,1]\) 上可导，\(f(0)=0, f(1)=2\) ．证明：在 \((0,1)\) 上存在两个 \(\xi_{1}\) 与 \(\xi_{2}\) ，使得

\[
\frac{1}{f^{\prime}\left(\xi_{1}\right)}+\frac{1}{f^{\prime}\left(\xi_{2}\right)}=1
\]

\subsection{题目4}
4．（15 分）证明：若函数 \(f(x)\) 在 \([a,+\infty)\) 上连续，且 \(\lim _{x \rightarrow+\infty}[b x-f(x)]=0\) ，其中 \(b\) 为非零实数．则\(f(x)\) 在 \([a,+\infty)\) 上一致连续。

\subsection{题目5}
5．（15 分）证明：存在常数 \(A<1\) ，对任意的正整数 \(n>1\) ，都有 \(\sqrt{1}+\sqrt{2}+\cdots+\sqrt{n}<A n^{\frac{3}{2}}\) ．

\subsection{题目6}
6．（15 分）判别函数列 \(\left\{n x(1-x)^{n}\right\}\) 在 \([0,1]\) 上的一致收敛性．

\subsection{题目7}
7．（15 分）求幂级数 \(\sum_{n=1}^{\infty} n^{2} x^{n+1}\) 的收敛域与和函数．

\subsection{题目8}
8．（15 分）求函数 \(f(x, y, z)=x-2 y+2 z\) 在条件 \(x^{2}+y^{2}+z^{2}=1\) 的约束下的最大值和最小值．

\subsection{题目9}
9．（15 分）求函数 \(f(x)=\int_{0}^{+\infty} \frac{\ln \left(1+y^{3}\right)}{y^{x}} \mathrm{~d} y\) 的定义域．

\subsection{题目10}
10．（15 分）记

\[
F(t)=\iiint_{V} f\left(x^{2}+y^{2}+z^{2}\right) \mathrm{d} x \mathrm{~d} y \mathrm{~d} z .
\]

其中 \(V=\left\{(x, y, z) \mid x^{2}+y^{2}+z^{2} \leq t^{2}\right\}(t>0)\) ，若 \(f\) 可微，求 \(F^{\prime}(t)\) ．
\end{document}
